This work introduces YANA, an open-source neuromorphic accelerator designed to bridge the simulation-to-hardware gap that currently limits the advancement of \ac{SNN} research and development. By providing accessible neuromorphic computing capabilities on affordable FPGA platforms, YANA enables researchers and practitioners to move beyond simulation-only SNN development and validate their algorithms on actual event-driven hardware.

YANA's architecture demonstrates several key contributions to the neuromorphic computing field. Its event-driven processing pipeline efficiently exploits both temporal and spatial sparsity inherent in SNNs, achieving near-linear scaling of inference time with sparsity levels.
Furthermore, YANA's support for arbitrary SNN topologies through point-to-point connections is future proof for the evolving landscape of neuromorphic applications, allowing for flexible network designs that can adapt to various research needs.

Our experimental evaluation on the Spiking Heidelberg Digits dataset validates YANA's ability to leverage sparsity for performance improvements, while the deployment on the accessible AMD Kria KR260 platform demonstrates practical resource efficiency with manageable hardware requirements. The YANA software framework, from training through deployment, provides a complete solution that integrates with the existing neuromorphic open-source ecosystem.
As the neuromorphic computing field continues to mature, accessible platforms like YANA will be crucial for translating the promising theoretical advances in SNNs into practical, real-world applications that fully exploit the unique advantages of brain-inspired computing paradigms.